\chapter{Electromagnetic Fields, Maxwell's Equations, and Self-Consistent Coupling}\label{ch:3}

\begin{chaptersummary}
The preceding chapters constructed matter models from particles and distributions. Electromagnetic fields now couple self-consistently to that matter: Maxwell constraints express charge conservation, field--particle feedback produces nonlinearity, and linearization extracts the operator governing small disturbances. This local response underpins dispersion, stability, and the electromagnetic discretizations of Volume II.
\end{chaptersummary}
\learninggoals{Connect charge continuity to Maxwell constraints; identify the data exchanged in Vlasov--Maxwell coupling; construct a Fr\'echet derivative at a specified base state and state the conditions for plane-wave analysis.}
\studyroute{Specific material models and kinetic responses are omitted in favor of the common structure of linearized operators.}{Fr\'echet derivative; linear response; Vlasov linearization}{Find the tangent tensor for $\bm D=\varepsilon\bm E+\alpha|\bm E|^2\bm E$ and explain how the base-state direction enters the linearized equations.}

\section{Electromagnetic Conservation and Field--Particle Coupling}\index{Vlasov--Maxwell system}
Maxwell's equations are
\[
\nabla\cdot\bm E=\frac{\rho_c}{\varepsilon_0},
\qquad
\nabla\cdot\bm B=0,
\]
\[
\frac{\partial\bm B}{\partial t}+\nabla\times\bm E=0,
\]
\[
\frac{\partial\bm E}{\partial t}
-c^2\nabla\times\bm B
=
-\frac{\bm J}{\varepsilon_0}.
\]

Taking the divergence of the Amp\`ere--Maxwell equation and using Gauss's law gives
\[
\frac{\partial\rho_c}{\partial t}
+\nabla\cdot\bm J=0.
\]

\begin{askai}
Ask AI:
\begin{quote}
Starting from the Amp\`ere--Maxwell equation
\[
\partial_t\bm E-c^2\nabla\times\bm B=-\bm J/\varepsilon_0
\]
take the divergence and use Gauss's law to derive charge conservation. Explain the structural role of \(\nabla\cdot(\nabla\times\bm B)=0\).
\end{quote}

Follow-up:
\begin{quote}
Construct discrete operators that do not satisfy \(D_hC_h=0\), and show how they can generate nonphysical charge.
\end{quote}
\end{askai}

Charge conservation is therefore built into the continuous structure of Maxwell's equations. A discretization that breaks this structure can generate nonphysical charge.

The Lorentz force density in a charged continuum is
\[
\bm f_L=\rho_c\bm E+\bm J\times\bm B.
\]
Combining fluid equations with Ohm's law gives two-fluid models or single-fluid MHD.
For kinetic species, $f_s$ determines $\rho_c$ and $\bm J$, while the fields change $f_s$ through $q_s(\bm E+\bm v\times\bm B)$. This two-way dependence is self-consistent field--particle coupling. Single-fluid MHD additionally uses low-frequency, quasineutral, and scale-ordering assumptions to reduce particle and species-fluid descriptions to a common velocity, pressure, and magnetic field. Chapter~\ref{ch:18} gives the closure.

\section{Nonlinear Operators and Local Linearization}
Vacuum Maxwell equations are linear. Coupling them to media, carriers, fluids, or nonlinear constitutive laws generally produces a nonlinear system. Write the coupled model abstractly as
\[
\Nop(\bm U)=0.
\]

\begin{authornote}
\textbf{Potential Question:}
"How does writing the coupled Maxwell system as \(\Nop(\bm U)=0\) help analyze perturbations?"

\textbf{Author Supplement:}
Time derivatives, spatial derivatives, and nonlinear couplings can be collected into an operator acting on the unknown fields. Without changing the problem, this notation provides a common description of perturbations near a solution. The operator derivative describes the propagation of small changes and supplies the linear equation approximately solved at each Newton step.

\textbf{Core Conceptual Connection:}
\[
\text{coupled field equations}
\longrightarrow
\Nop(\bm U)=0
\longrightarrow
\text{linearized operator}
\]
\end{authornote}

Near a base state \(\bm U_0\), write
\[
\bm U=\bm U_0+\epsilon\bm u,
\qquad 0<\epsilon\ll1.
\]
For a sufficiently smooth operator,
\[
\Nop(\bm U_0+\epsilon\bm u)
=
\Nop(\bm U_0)
+\epsilon D\Nop[\bm U_0](\bm u)
+\frac{\epsilon^2}{2}
D^2\Nop[\bm U_0](\bm u,\bm u)
+\cdots.
\]
If the base state satisfies
\[
\Nop(\bm U_0)=0,
\]
retaining first-order terms gives
\[
\Lop\bm u=0,
\qquad
\Lop:=D\Nop[\bm U_0].
\]

Here \(D\Nop[\bm U_0]\) is the Fr\'echet derivative, also called the linearized, Jacobian, or tangent operator. It is the local linear approximation of the nonlinear operator at the base state.

For example, let a nonlinear medium satisfy
\[
\bm D=\bm D(\bm E),
\qquad
\bm B=\bm B(\bm H).
\]
Set
\[
\bm E=\bm E_0+\epsilon\bm e,
\qquad
\bm H=\bm H_0+\epsilon\bm h.
\]
The first-order constitutive expansion is
\[
\bm D(\bm E_0+\epsilon\bm e)
=
\bm D(\bm E_0)
+
\epsilon
\left.
\frac{\partial\bm D}{\partial\bm E}
\right|_{\bm E_0}\bm e
+\Order(\epsilon^2),
\]
\[
\bm B(\bm H_0+\epsilon\bm h)
=
\bm B(\bm H_0)
+
\epsilon
\left.
\frac{\partial\bm B}{\partial\bm H}
\right|_{\bm H_0}\bm h
+\Order(\epsilon^2).
\]

\begin{askai}
Ask AI:
\begin{quote}
Compare finite-dimensional Taylor expansion with the Fr\'echet derivative, and explain step by step why the nonlinear material relation
\[
\bm D=\bm D(\bm E),\qquad \bm B=\bm B(\bm H)
\]
produces a tangent material tensor near a base state. Use an anisotropic matrix example.
\end{quote}

Follow-up:
\begin{quote}
If \(\bm D(\bm E)\) is nonlinear, why does the perturbation equation contain \(\partial\bm D/\partial\bm E|_{\bm E_0}\) rather than the original \(\bm D\)?
\end{quote}
\end{askai}

Define the tangent material tensors
\[
\bm\varepsilon_{\mathrm{tan}}
=
\left.\frac{\partial\bm D}{\partial\bm E}\right|_{\bm E_0},
\qquad
\bm\mu_{\mathrm{tan}}
=
\left.\frac{\partial\bm B}{\partial\bm H}\right|_{\bm H_0}.
\]
For instantaneous, nondispersive constitutive laws, further assume a time-independent base state, so the tangent tensors are constant in time, and set the perturbation current to zero. The first-order perturbations then satisfy a linear Maxwell-type system:
\[
\nabla\times\bm e
=
-\bm\mu_{\mathrm{tan}}\frac{\partial\bm h}{\partial t},
\]
\[
\nabla\times\bm h
=
\bm\varepsilon_{\mathrm{tan}}\frac{\partial\bm e}{\partial t}.
\]

\section{Vlasov Perturbations and Wave Dispersion}
In the Vlasov--Maxwell system, choose an equilibrium $f_{s0}$, $\bm E_0=0$, and specified $\bm B_0$ satisfying the zeroth-order equations, and write
\[
f_s=f_{s0}+\epsilon f_{s1},
\qquad
\bm E=\epsilon\bm E_1,
\qquad
\bm B=\bm B_0+\epsilon\bm B_1.
\]
Neglecting \(\Order(\epsilon^2)\) terms gives
\[
\frac{\partial f_{s1}}{\partial t}
+\bm v\cdot\nabla_{\bm x}f_{s1}
+\frac{q_s}{m_s}
(\bm v\times\bm B_0)\cdot\nabla_{\bm v}f_{s1}
=
-\frac{q_s}{m_s}
(\bm E_1+\bm v\times\bm B_1)
\cdot\nabla_{\bm v}f_{s0}.
\]
The perturbation current is
\[
\bm J_1=\sum_s q_s\int \bm v f_{s1}\dd\bm v.
\]

For a closed constant-coefficient linear system with finitely many amplitude unknowns, or a kinetic system whose velocity response has already been integrated, use plane waves:
\[
\bm u_1(\bm x,t)
=
\widehat{\bm u}
\exp[\ii(\bm k\cdot\bm x-\omega t)],
\]

\begin{authornote}
\textbf{Potential Question:}
"Why use plane waves to analyze perturbations? Does this assume that every solution is a wave?"

\textbf{Author Supplement:}
A constant-coefficient operator linearized about a uniform background preserves each Fourier wave number, allowing separate modal analysis. Slowly varying coefficients can be frozen locally to probe local response with plane waves. Boundaries and global nonuniformity enter through the corresponding global problem.

\textbf{Core Conceptual Connection:}
\[
\text{linear perturbation equation}
\longrightarrow
\text{plane-wave ansatz}
\longrightarrow
\text{dispersion relation}
\]
\end{authornote}

Then
\[
\partial_t\mapsto-\ii\omega,
\qquad
\nabla\mapsto\ii\bm k.
\]
and the linear PDE system becomes
\[
\bm M(\omega,\bm k)\widehat{\bm u}=0.
\]

\begin{askai}
Ask AI:
\begin{quote}
Starting with the plane-wave substitution
\[
\partial_t\mapsto-\ii\omega,\qquad \nabla\mapsto\ii\bm k
\]
choose a simple linear PDE system and show step by step how it becomes
\[
\bm M(\omega,\bm k)\widehat{\bm u}=0.
\]
Explain why a nonzero amplitude requires \(\det\bm M(\omega,\bm k)=0\).
\end{quote}

Follow-up:
\begin{quote}
Derive the complete dispersion relation for one-dimensional acoustic waves or vacuum Maxwell plane waves.
\end{quote}
\end{askai}

The condition for a nonzero solution,
\[
\det\bm M(\omega,\bm k)=0
\]
is the dispersion relation.

For kinetic equations that retain continuous velocity variables, the first result is a velocity-space response problem. An analytic continuation, integration contour, and field equations must be specified before the corresponding dielectric response and dispersion condition can be constructed.

Perturbation theory constructs solutions order by order in a small parameter; stability analysis examines linear perturbation growth rates. The Fr\'echet derivative also generates the Jacobian used in nonlinear algebraic iteration, whose solution algorithms are developed in Volume~0.
